\section{Discussion}
\label{sec:discussion}

\subsection{Interpretation of the performance evidence}

Four consequences follow, each contradicting a plausible default assumption. The recorded authorization-associated span accounted for roughly a quarter of the historical write path at 4 peers and varied less across the tested peer configurations than the uninstrumented remainder, rising to \SpanSharePeerThirtyTwo\% of the smaller total at 32 peers. Because its historical instrumentation boundaries are undocumented, these shares are descriptive and do not establish a transferable optimisation ceiling (Limitation L1). Because expired-role cleanup mutates state inside \texttt{CheckAccess}, read-intent checks are endorsed rather than served as peer queries; a read-only decision path with expiry reclaimed by a background sweep would take reads off the consensus path, at the cost of the durable ordered audit transaction a committed \texttt{CheckAccess} produces and needing revocation freshness maintained another way, so we call it a high-value candidate worth that trade-off in most deployments rather than an unconditional recommendation. Peers are associated with lower tail latency here; treat throughput separately. Configured logical peer multiplicity is associated with a close-to-log-linear reduction in write latency (Section~\ref{sec:results-peer-scaling}), supported by a joint regression on peer count and run order once the earlier order confound is designed out (L2), while the separate, interleaved concurrency experiment found the same observed throughput peak for both arms without locating the saturating resource (L3): ``peers buy latency, not capacity'' describes only the latency side, and a deployment weighing more endorsers against more capacity should measure both under its own workload. Budget the complete access-control-enabled path empirically. Across the tested closed-loop concurrency levels, this path averaged \ThroughputPctDiffMean\% lower throughput than the comparison path, and no condition-by-concurrency interaction was detected. For the tested hardware, peer configuration, transaction semantics and closed-loop workload, the observed mean difference provides a provisional planning estimate; it should be re-estimated for other deployments and should not be interpreted as a universal authorization derating factor. Treat revocation as an unresolved operational risk. \RevocationDelayedCount{} of \RevocationStageRecords{} stage records carry the deployment's delayed label; their association with recorded connectivity outages motivates a reachability hypothesis, but the trace cannot determine per-episode cause or exclude the propagation protocol as a contributor (L5); event-driven credential status checking \citep{w3c2024bitstring} and constrained-device status protocols such as TinyOCSP \citep{fedrecheski2023tinyocsp} remove the dependency on a scheduled publication reaching a possibly-disconnected gateway, whichever the cause turns out to be.

Direct numerical comparison against prior systems is limited by methodological heterogeneity: hardware, workload, transaction semantics and what is counted as latency all differ (supplementary material), a point Abang et al.\ \citep{abang2024latency} make about the Fabric latency literature generally; the comparison we regard as meaningful is the internal one, the same network, hardware and workload with and without the access-control layer (Section~\ref{sec:results-throughput}), since it is reproducible from the released artifact.

For a practitioner deciding whether to adopt an HRBAC layer of this shape on similar edge hardware, three operational recommendations follow directly from what was actually measured, each qualified by the limitation that licenses it. First, budget the authorization-associated span as a roughly fixed cost near \SpanLatencyPeerFour{} ms rather than a percentage of the write path, since Section~\ref{sec:span-variation} shows the percentage moves with peer count for reasons unrelated to the access-control logic itself (L1). Second, treat additional logical endorsers as a latency lever, not a throughput lever, and measure both independently before provisioning a deployment on their basis (L2, L3). Third, treat certificate revocation as an operational process requiring its own monitoring rather than a mechanism whose correctness the deployment traces already demonstrate; the observability gap in Section~\ref{sec:revocation-audit} means an operator cannot currently answer "how long was this credential live after revocation" from the released fields alone (L5). None of these recommendations transfers to the specified-but-unimplemented corrected artifacts (L6); they describe operating the system as it was actually deployed.

\subsection{Cross-layer lessons from the deployment postmortem}

The three defect classes went undetected because nothing compared the implementation against its description: the policy diagram and the inheritance table agreed with each other and disagreed with the code, the encoder guarded a bound the producer violated, and the verifier and signer used different message formats, while tests for each component in isolation passed. Reverifying these findings directly against the released source for this version, rather than trusting the earlier draft's prose, changed none of the qualitative conclusions and confirmed each defect's exact location (file and function) in Sections~\ref{sec:audit-policy}--\ref{sec:revocation-audit}. None of the three defects is sophisticated; their absence of detection allowed a \DeploymentDays{}-day deployment to operate with an implementation that contradicted its own specification.

The broader lesson is methodological. Performance measurements can remain internally valid while the security semantics of the measured system are wrong. A deployment postmortem should therefore pair trace analysis with code-to-specification review and cross-component tests, and should state explicitly which historical variables remain descriptive, require reinterpretation or must be withdrawn (Table~\ref{tab:validity-table}). A second, narrower lesson follows from Section~\ref{sec:corrected-absent}: describing a correction in prose is not the same as releasing it, and a reanalysis that only re-derives numbers from data without re-verifying which artifacts actually exist in the repository can silently inherit a documentation gap as if it were a shipped fix. This version treats that check as part of the audit, not an afterthought.

\subsection{Security and validity boundaries}

The threat model treats orderers and the certificate authority as non-malicious under crash-fault-tolerant Raft, and physical tampering and side-channel extraction as residual risk, architecturally addressed by mutual TLS and constant-time primitives rather than empirically tested. Within that scope, three things this deployment measured are not the same as validating the intended design. Section~\ref{sec:results-security} measures refusals against the deployed policy, which Section~\ref{sec:audit-policy} shows is not the intended one, and its role/operation oracle is our own disclosed construction rather than a shipped independent artifact (L4), so it is conformance testing of one implementation under one analyst-supplied mapping, not a security proof, and says nothing about an adversary who deviates from the scripted campaign. The postmortem's corrections are specifications, not patches to running code: as Section~\ref{sec:corrected-absent} states plainly, no corrected chaincode, gateway verifier or firmware exists in this repository, so no corrected-implementation performance, correctness or security claim is made (L6). And overhead throughout is attributed by differencing against a baseline with the access-control chaincode removed, which also removes the audit write and the nonce check, so every reported cost bounds the access-control layer rather than an isolated permission decision. Table~\ref{tab:validity-table} states, trace by trace, what remains usable after the audit.

\subsection{Limitations and required validation}
\label{sec:limitations}

This study has six principal limitations. Table~\ref{tab:limitations} separates each unresolved measurement boundary from its consequence for interpretation. The evidence needed to reduce each limitation is described alongside it in the supplementary material.

\begin{center}
\small
\begin{longtable}{p{2.0cm}p{5.0cm}p{5.6cm}}
\caption{Principal limitations and their consequences for interpretation.}
\label{tab:limitations}\\
\toprule
Limitation & Unresolved boundary & Consequence for interpretation \\
\midrule
\endfirsthead
\multicolumn{3}{l}{\small\itshape Table~\ref{tab:limitations} continued} \\
\toprule
Limitation & Unresolved boundary & Consequence for interpretation \\
\midrule
\endhead
\bottomrule
\endfoot
L1 -- Timer boundaries & The exact source-code boundaries of the historical \texttt{rbac\_overhead\_ms} timer cannot be reconstructed. The pooled operation and role labels also do not isolate hierarchy depth, privilege level or denial mechanism. & The value is an implementation-defined authorization-associated span, not an isolated policy-evaluation cost. Stage attribution requires timestamps at endorsement, ordering, validation and commit boundaries. \\
L2 -- Peer multiplicity and external validity & The peer experiment varied logical peers on four fixed physical hosts (single \texttt{host\_id} for all runs). Counterbalancing controls execution position only within this campaign, and throughput was not measured as a function of peer count. & The experiment does not establish physical scale-out effects or generalisation to other hardware, sites, seasons or workloads, and it cannot support a peer-count throughput claim. \\
L3 -- Throughput campaign order and saturation & The two conditions were interleaved within each concurrency level, but the levels were tested in one fixed ascending sequence within a single session. Concurrency was sampled in ten-client steps, and the saturated resource was not instrumented. & The within-level position term cannot exclude campaign-wide drift confounded with concurrency. The true peak may lie between tested levels, and the experiment does not identify the resource responsible for saturation. \\
L4 -- Authorization-corpus coverage & The historical corpus contains denied requests but no permitted requests. Role--operation outcomes were checked against an oracle we disclose and construct for this version (no \texttt{policy-requirements.json} exists in the repository); zone, temporal, replay and identity dimensions remain scored against \texttt{roles.go}-derived logic only. & The corpus cannot estimate a false-denial rate or establish complete conformance to the intended policy; the confirmed/discrepant split is sensitive to our disclosed operation-to-permission mapping. \\
L5 -- Revocation observability & Revocation records are lifecycle stages linked one-to-one to subjects in this trace, but a subject revoked more than once over the deployment cannot be disambiguated, and no stage-to-stage episode key is carried through the pipeline. & End-to-end exposure and causal attribution cannot be reconstructed from the historical trace. \\
L6 -- No corrected implementation in this repository & No corrected chaincode, firmware or gateway package exists in the repository inventoried for this version, despite the corrections being fully specified and arithmetically checked (Section~\ref{sec:corrected-absent}). Historical energy was measured on one device, and the deployment covered one site, one season and a sensor-write-dominated workload. & Historical latency, throughput, cryptographic and energy results describe the deployed implementation and must not be assumed to transfer to any future corrected version, because none has been built, tested or measured. Numerical estimates are deployment-specific rather than universal IoT or Fabric constants. \\
\end{longtable}
\end{center}
